\section{Defining the Task}
\label{sec:defining-task}

RP began as a bachelor task offered by Admmit, a Norwegian software firm building products for the transport sector. The task description, distributed through NTNU as project OPG29, set out a brief: investigate how Norwegian traffic coordinators plan daily transport assignments, and build a working prototype of a planning platform that could assist them. The team accepted the task and committed to delivering both the artefact and the thesis that surrounds it. From that decision onwards, two roles ran in parallel for the two-student team: building RP as a functioning platform, and conducting the research that would explain and justify its design choices. Both students contributed to both workstreams over the course of the project.

The brief defined the topic but not the empirical foundation, and that part of the work was arranged by the team rather than by Admmit. Seven traffic coordinators and transport managers were contacted directly by phone over a single day in March 2026. All seven worked at companies that were already Admmit customers, but Admmit did not broker the introductions. The team found the coordinators through Admmit's customer roster, called them on its own initiative, and conducted brief semi-structured interviews lasting roughly fifteen minutes each. The interviewed companies range from a small operator with eight to twenty drivers to a multi-department fleet of around forty-five vehicles, and they use a mix of Timpex, Opter, internal tools, and fully manual planning.

Before any coordinator was on the phone, however, two design choices had already been locked. Admmit specified Human-in-the-Loop (HITL) operation as a non-negotiable property of the artefact from the project's first week: the algorithm proposes a plan; the coordinator inspects, modifies, accepts, or rejects each assignment. Admmit also specified a multi-tenant architecture in which a single deployed instance serves multiple customer companies through configuration. Both decisions originate in Admmit's mandate, not in the interviews. The interviews validated HITL by surfacing the coordinator's insistence on retaining authority over assignments, and they validated the cross-company variety that motivates configurable behaviour, but they did not introduce either property. The origin map presented later in this chapter (\Cref{sec:iterations}) labels every major design decision by origin and by interview-validation status, so the reader can tell apart Admmit-mandated choices, interview-driven choices, and designer / domain-knowledge choices.

This origin frames the rest of the methodology. Every design choice traced through the chapters that follow comes from one of three sources: the Admmit mandate that fixed HITL and multi-tenancy at the outset, the dialogue with the seven coordinators that informed the planning interface and the assignment-criteria taxonomy, or designer and domain-knowledge decisions that filled in what neither the mandate nor the interviews specified. HITL turns the \textbf{Trust/control} anchor into four concrete coordinator actions on every algorithm-generated assignment. The multi-tenant architecture supports \textbf{Adaptability} by isolating per-customer configuration in a single deployed instance. The interview-driven analysis identifies the \textbf{Efficiency} anchor as the resource-utilization visibility gap that the artefact's planning interface and benchmark are designed to address. The remaining sections of this chapter describe the research paradigm under which the artefact was built (\Cref{sec:dsr-dsrm}) and how the interview and analysis work was conducted (\Cref{sec:data-collection,sec:data-analysis}). The eight named iterations that produced the system are documented in \Cref{sec:iterations}. The framework under which it was evaluated, together with the validity bounds that follow, is set out in \Cref{sec:evaluation-framework,sec:validity}.

